\section{System model and problem formulation}
\label{sec:prelim}

\paragraph{System model.}
Let $\clientset = \{1, \dots, N\}$ denote the client set, and let each client $i \in
\clientset$ hold a local data stream $D_i = (D_i^{(1)}, D_i^{(2)}, \dots)$, where
$D_i^{(t)}$ is the data available to client $i$ at round $t$. At each round $t$ of an
FCL protocol, the server broadcasts a global model $\model^{(t)}$; a subset $S^{(t)}
\subseteq \clientset$ of clients trains locally on $D_i^{(t)}$ and returns an update
$\Delta_i^{(t)}$; the server aggregates $\{\Delta_i^{(t)}\}_{i \in S^{(t)}}$ to produce
$\model^{(t+1)}$. We denote the aggregation operator $\mathcal{A}$ and the full training
trajectory as $\Theta_{r} = (\model^{(0)}, \dots, \model^{(r)})$. An \emph{erasure
request} is a tuple $(\forget, r_{\forget}, r_{\mathrm{req}})$: client $\forget \in
\clientset$ participated up to round $r_{\forget}$ and issues the request at round
$r_{\mathrm{req}} \geq r_{\forget}$. We denote by $\Theta_{-\forget, r}$ the
counterfactual training trajectory in which client $\forget$ never participated, i.e., a
full retrain from scratch without their data. The gold standard is to produce a model
$\model^{(r)}$ distributed indistinguishably from $\model_{-\forget}^{(r)}$. For the
sake of readability, an overview of the notations used throughout this paper is provided
in Table~\ref{tab:notations}.

\begin{table}[t]
\centering
\scriptsize
\setlength{\tabcolsep}{3pt}
\begin{tabular}{ll}
\toprule
Symbol & Meaning \\
\midrule
$\clientset,\, N$              & Client set and its size \\
$D_i^{(t)}$                    & Local data of client $i$ at round $t$ \\
$S^{(t)}$                      & Clients sampled at round $t$ \\
$\model^{(t)} \!=\! (\theta_\mathrm{sh}, \theta_\mathrm{cl})$ & Global model: features / head \\
$\Delta_i^{(t)}$               & Local update of client $i$ at round $t$ \\
$\mathcal{A}$                  & Server aggregation operator \\
$\Theta_r$                     & Training trajectory up to round $r$ \\
$\forget$                      & Departing client's index \\
$r_\forget,\, r_\mathrm{req}$  & Departure / request rounds \\
$T$                            & Cooling-off period $r_\mathrm{req} - r_\forget$ \\
$\Theta_{-\forget, r}$         & Counterfactual trajectory w/o $\forget$ \\
$C,\, \alpha,\, \sigma$        & Clip norm, decay, server noise \\
$\sigma_i$                     & Pre-committed DP summary of client $i$ \\
$(\eps_\summary, \delta_\summary)$ & Per-summary DP budget \\
$\mu_1(T)$                     & GDP parameter of feature-block release \\
$\mu_\summary$                 & GDP parameter of head-subtraction release \\
$\mu_\mathrm{joint}(T)$        & GDP parameter of composed joint release \\
$(\eps_\mathrm{total}, \delta_\mathrm{total})$ & Overall formal guarantee \\
$\gamma,\, R$                  & Ascent step / projection radius \\
$\rho$                         & Client participation rate \\
$M(X) \indist M(X')$           & $(\eps, \delta)$-indistinguishability \\
\bottomrule
\end{tabular}
\caption{Notation.}
\label{tab:notations}
\end{table}

\paragraph{Differential privacy.}
A randomised mechanism $M: \mathcal{X} \to \mathcal{Y}$ is $(\varepsilon,
\delta)$-differentially private if, for all neighbouring datasets $X, X'$ differing in a
single record (or, at client level, a single client), and all measurable $Y \subseteq
\mathcal{Y}$,
\begin{equation}
\Prob[M(X) \in Y] \leq e^{\varepsilon} \Prob[M(X') \in Y] + \delta.
\label{eq:dp}
\end{equation}
We write $M(X) \indist M(X')$ when \eqref{eq:dp} holds. Differential Privacy (DP) is
closed under post-processing: for any (possibly randomised) $f$, $f \circ M$ is also
$(\varepsilon, \delta)$-DP.

\paragraph{Threat model.}
We assume an honest-but-curious server that follows the protocol but may attempt to
infer private information from received messages. This is stronger than the
malicious-server assumption of some FU work but matches the deployed setting for most FL
systems. An external adversary observes any released model artefact (before or after
erasure) and may query it as a black box. This adversary implements MIA
\cite{shokri2017mia,carlini2022lira} and FUIA \cite{fuia2025}. Clients are
semi-honest: they follow the training protocol but may attempt to influence the erasure
procedure adversarially (e.g., contribute a backdoor before requesting erasure).
Compromised transmissions are out of scope; we assume standard secure aggregation and
TLS between clients and server. The controller must be able to produce evidence,
examinable by a supervisory authority, that after erasure the model is $(\varepsilon,
\delta)$-indistinguishable from a counterfactual retrain, without publishing a
pre-erasure model to any party. The supervisory authority is modelled as a semi-honest
protocol participant, not an external third party: it receives read access to the
pre-erasure head weights under a non-disclosure obligation in order to recompute the
head delta during verification (Section~\ref{sec:method-erasure}). The Chen et al.\
unlearning-leakage channel \cite{chen2021unlearningleaks} is therefore closed with
respect to parties outside the controller–authority pair, which includes the
departing client, the remaining clients, and any external observer.

The scope of the formal guarantee is restricted to the departing client. Specifically,
FORGE provides a client-level $(\varepsilon, \delta)$-indistinguishability bound against
the counterfactual trajectory in which client $\forget$ never participated. Privacy for
non-departing clients inherits from the underlying FL protocol (secure aggregation, and
DP-FedAvg \cite{mcmahan2018learning} if deployed in combination). The composition
theorem (Theorem~\ref{thm:composition}) does not extend non-departing-client bounds
beyond the base-protocol level. This matches the scope of all prior FU work that we
compare against.

An adversarial departing client that recorded their per-round received broadcasts can,
by assumption, reconstruct a pre-erasure global model at the round of their departure.
In such a scenario, the Chen-et-al channel \cite{chen2021unlearningleaks} is closed on
the server side as follows: the erasure certificate issued at the end of
Algorithm~\ref{alg:erasure} comprises a hash of the pre-committed summary and a
proof-of-execution transcript (defined in Section~\ref{sec:method-erasure}), neither of
which exposes the post-erasure model. The departing client may therefore verify that
an erasure took place without observing the post-erasure model artefact, and no third
party receives a before/after pair.

\paragraph{Forward forgetting.}
We now introduce the central definition of this paper.

\begin{definition}[$(\varepsilon, \delta, T)$-forward forgetting]
\label{def:forward-forgetting}
An FCL protocol satisfies \emph{$(\varepsilon, \delta, T)$-forward forgetting} for
client $\forget$ if, for any erasure request $(\forget, r_{\forget}, r_{\mathrm{req}})$
with $r_{\mathrm{req}} \geq r_{\forget} + T$, the released model
$\model^{(r_{\mathrm{req}})}$ satisfies
\begin{equation}
\model^{(r_{\mathrm{req}})} \indist \model_{-\forget}^{(r_{\mathrm{req}})}.
\end{equation}
The protocol satisfies \emph{$(\varepsilon(T), \delta(T))$-forward forgetting} if the
guarantee holds for all $T \geq 0$ with rate functions $\varepsilon(T), \delta(T)$.
\end{definition}

Forward forgetting differs from standard DP in three ways. First, the reference
distribution is the counterfactual retrain trajectory, not a database-neighbours pair.
Second, the guarantee is \emph{time-indexed}: the $\varepsilon$ is a function of how
long the client has been absent, not a single number set at training. Third, the
guarantee is written in a one-sided form, i.e., it bounds the distance from the
released model to the counterfactual.

\begin{remark}[Equivalence to two-sided DP under symmetric sampling]
\label{rem:two-sided}
Under exchangeable client sampling $S^{(t)}$ and i.i.d.\ local-training randomness
(Assumption~\ref{ass:noise}), the forward and reverse directions
$(\model^{(r_\mathrm{req})}, \model_{-\forget}^{(r_\mathrm{req})})$ are equidistributed
with respect to swapping client $\forget$ with any other sampled client. The one-sided
bound of Definition~\ref{def:forward-forgetting} is therefore equivalent to the standard
two-sided $(\varepsilon, \delta)$-DP neighbouring-datasets condition up to a factor of
at most $2$ in $\delta$, by the group-privacy argument of \cite{dwork2014book}, Section
2.3.
\end{remark}

\begin{remark}
Definition~\ref{def:forward-forgetting} does not require the erasure request to trigger
any procedure at all. If the protocol's forward-forgetting rate $\varepsilon(T)$ is
small enough at $T = r_{\mathrm{req}} - r_{\forget}$, the released model already
satisfies the required indistinguishability. In practice we combine decay-based forward
forgetting (which handles the shared-signal component and covers the bulk of the
guarantee for $T \geq T^{\star}$) with an active subtraction step (which covers the
memorised residual regardless of $T$).
\end{remark}

\paragraph{Assumptions.}
We make the following standard assumptions.

\begin{assumption}[Bounded update norm]
\label{ass:clipping}
Each client update satisfies $\|\Delta_i^{(t)}\| \leq C$ for a clipping constant $C > 0$.
This is enforced by clipping at the client.
\end{assumption}

\begin{assumption}[Lipschitz loss]
\label{ass:lipschitz}
The local loss $\ell(\model; x)$ is $L$-Lipschitz in $\model$ for all $x$.
\end{assumption}

\begin{assumption}[Bounded participation]
\label{ass:participation}
Each client participates in at most $K$ rounds out of any window of $T$ consecutive
rounds, with $K/T \leq \rho$ for a participation rate $\rho \in (0, 1)$.
\end{assumption}

\begin{assumption}[Sub-Gaussian gradient noise]
\label{ass:noise}
Local minibatch gradient noise is sub-Gaussian with parameter $\eta^2$.
\end{assumption}

Assumptions~\ref{ass:clipping}--\ref{ass:noise} are standard for federated DP analysis
\cite{mcmahan2018learning,abadi2016dpsgd}. Assumption~\ref{ass:participation} bounds the
frequency with which any single client can inject signal, and directly controls the
decay rate.

\paragraph{Problem statement.}
Given the FCL setup above, we aim to design a protocol satisfying the following four
properties.
\begin{itemize}[leftmargin=*,itemsep=1pt]
\item[\textbf{P1.}] \emph{Forward forgetting}: $(\varepsilon(T), \delta(T))$-forward
forgetting for every client, with $\varepsilon(T) \to 0$ as $T \to \infty$.
\item[\textbf{P2.}] \emph{No history storage}: the server retains no per-round client
updates. Only the current model $\model^{(t)}$, aggregated statistics, and a hash of
each client's pre-committed summary may persist across rounds.
\item[\textbf{P3.}] \emph{No client cooperation at request time}: the erasure procedure
at round $r_{\mathrm{req}}$ does not require any communication with client $\forget$
beyond receiving the request itself.
\item[\textbf{P4.}] \emph{Bounded utility cost}: expected test accuracy of the final
released model degrades by at most $\Delta_{\mathrm{util}}$ compared to a non-private
FedAvg baseline, for a specified budget $\Delta_{\mathrm{util}}$.
\end{itemize}

We show in the remainder of the paper that no prior FU method satisfies P1--P3
simultaneously with a non-trivial $\Delta_{\mathrm{util}}$, and we present a protocol
that does.
